\documentclass[11pt]{article}
\usepackage{amsmath,amssymb,booktabs,graphicx,geometry}
\geometry{margin=0.85in}
\setlength{\tabcolsep}{4pt}

\title{Pathway attribution with TabDDPM and a meta-learner}
\author{}
\date{}

\begin{document}
\maketitle

\section{Pathway}

The comparison is not one batch against another. The anchor is a peak day.
The next seven days are seven forward steps. On each step a TabDDPM fit only
on the peak imputes one feature block, conditional on the other columns of
that day's row. The diffusion class is
\texttt{GaussianMultinomialDiffusion} from \texttt{yandex-research/tab-ddpm}
(commit \texttt{b476257}). Numerical columns use the Gaussian head. Known
columns are reset to the anchor noising process at every reverse step, the
RePaint conditioning step. The label passed to the denoiser is drawn from the
peak-day label distribution. Observed rows are clipped to \([-5,5]\) before
training, and the predicted clean sample is clamped to the same box. On
these tables the unclipped reverse chain leaves that box.

Predictive loss is the Brier score of a logistic model fit on the peak and
transferred onto day \(t\). On the synthetic panel and MovieLens a ridge on
the peak also reports rating or latent MSE. PO-risk follows PermuCATE on
rows pooled from the peak (\(W=0\)) and day \(t\) (\(W=1\)): random forests
for \(\mu_0\), \(\mu_1\), and \(\hat\tau\), logistic propensity with
clipping, and
\[
\phi(Z)=\frac{(Y-\mu_T)(T-\pi)}{\pi(1-\pi)}+(\mu_1-\mu_0),\qquad
\text{PO-risk}=\mathbb{E}\bigl[(\phi-\hat\tau(X))^2\bigr]
\]
evaluated on day \(t\) with \(W=1\). Replacing block \(b\) by its imputed
draw changes each loss by \(\Delta_b\). The rank share is
\(|\Delta_b|/\sum_{b'}|\Delta_{b'}|\). MMD compares the block on day \(t\)
with the same block on the peak. HSIC compares the block with the day's
label.

To separate sampler bias from attribution, Table~\ref{tab:sampler-floor}
reports a peak-day floor (redraw the block on the anchor) and the excess on
the last batch. A verdict of \texttt{sampler} means the later batch does not
move beyond that redraw; \texttt{drift} and \texttt{model} require a larger
excess, with \texttt{drift} also requiring a significant block MMD.

The last day reports a bootstrap percentile interval for each \(\Delta\) and
a 40-permutation \(p\)-value for MMD and HSIC. Across the seven steps, a
sign-flip test asks whether the block with the largest mean absolute
movement exceeds each other block.

\section{Synthetic peak}

Day 0 is standard normal, with four blocks of four columns: \texttt{stable},
\texttt{context}, \texttt{driver}, and \texttt{noise}. The outcome uses
\texttt{stable} and \texttt{driver}. Over the next seven days the
\texttt{driver} mean moves by up to \(1.5\) and its outcome coefficient
rises, while \texttt{context} receives a smaller location shift and no
outcome role. The unconditional TabDDPM sample has MMD \(0.000\) against
the peak day.

PO-risk rises from \(0.005\) to \(0.050\) and latent MSE from \(0.337\) to
\(4.496\). By day 7 the \texttt{driver} share is \(0.905\) of the PO-risk
movement, \(0.959\) of the Brier movement, \(0.936\) of the MSE movement,
and \(0.834\) of the MMD. Its PO interval is \([-0.122,-0.062]\) and its
Brier interval is \([-0.139,-0.126]\); both exclude zero. The sign-flip
\(p\)-values of \texttt{driver} against the other three blocks are at most
\(0.021\) for PO-risk, Brier, and MSE. Context has MMD \(p=0.024\), and its
Brier movement stays near \(0.001\). HSIC deltas are all below \(0.001\)
and do not rank the blocks.
Table~\ref{tab:pathway-synthetic}--\ref{tab:pathway-synthetic-sign} record
the path, the peak day, and the sign-flip tests.

\section{IEEE-CIS week after the fraud-rate peak}

The local fraud-rate peak with seven usable days after it is transaction
day 173 (fraud rate \(0.0495\)). The pathway runs through day 180
(\(0.0374\)). Six blocks are used: amount, card, address, counts,
timedelta, and hour. The monotone clock is not a feature. The unconditional
sample has MMD \(0.002\) against the peak day.

No block has an MMD \(p\)-value below \(0.05\) on day 180. The largest mean
absolute PO-risk movement is timedelta, and its sign-flip \(p\)-values
against card, address, and counts are all above \(0.178\). Every day-180 PO
interval includes zero. The transferred Brier movement is largest for card.
Its sign-flip \(p\)-value is \(0.008\) against amount and address and
\(0.028\) against counts; it is \(0.054\) against hour and \(0.292\) against
timedelta. Card's day-180 Brier interval is \([-0.023,-0.004]\), and its
MMD \(p\)-value is \(1\). Over this week the card block moves the fraud
model's predictions, and the covariate blocks do not separate.
Tables~\ref{tab:pathway-ieee-day} and~\ref{tab:pathway-ieee-sign} give the
day-180 numbers and the seven-step tests.

\section{MovieLens week after the high-rating peak}

The local peak in the rate of ratings of at least 4, among days that still
have seven later calendar days, is 2001-01-01 (rate \(0.615\)). The pathway
runs through 2001-01-08 (\(0.483\)). Rating is the outcome. The user block
is gender, age, occupation, and the first zip digit. The item block is
eight genre indicators. Context is the hour; weekday is constant inside one
calendar day and drops out. The peak day has 633 rows after keeping users
and items with at least 20 ratings. The unconditional sample has MMD
\(0.004\) against that day.

Mean MMD over the seven days is \(0.108\) for hour, \(0.100\) for user, and
\(0.012\) for genre. Hour against genre has sign-flip \(p=0.016\); hour
against user has \(p=0.413\). On the last day every block's MMD \(p\)-value
sits at the permutation floor \(0.024\), while the genre MMD itself is
\(0.006\).

Predictive loss points at genre. Mean absolute Brier movement is \(0.010\)
for genre, \(0.005\) for hour, and \(0.004\) for user. The sign-flip
\(p\)-value of genre against user is \(0.029\), and against hour \(0.074\).
On 2001-01-08 the genre Brier interval is \([-0.025,-0.010]\) and the user
Brier interval includes zero. Rating MSE also leads on genre (\(0.044\)
against \(0.027\) for user), with sign-flip \(p=0.125\) against user.
Mean absolute PO-risk movement is \(0.009\) for user and \(0.005\) for
genre, with sign-flip \(p=0.198\). The last-day user PO interval is
\([-0.048,-0.008]\); the genre PO interval includes zero. HSIC deltas stay
at most \(0.002\).

The covariate shift over this week is who is rating and at which hour.
The genre margin barely moves, and Brier still ranks genre first because
the peak-day rating model uses genre.
Tables~\ref{tab:movielens-path}--\ref{tab:movielens-sign} record the path.

\section{Home Credit along application id}

\texttt{application\_train.csv} has no event time. Rows sorted by
\texttt{SK\_ID\_CURR} are cut into 40 equal slices. The local peak in the
default rate is slice 31 (\(0.083\)), and the pathway runs to slice 38
(\(0.080\)). The five blocks are external (the three \texttt{EXT\_SOURCE} columns),
amount, personal, region, and other. Each slice contributes 700 rows. The
unconditional sample has MMD \(0.003\) against the peak slice.

Every block MMD stays below \(0.005\). The largest seven-step MMD is
external, and its sign-flip \(p\)-value against each other block is at
least \(0.266\). On the last slice only \texttt{other} has MMD \(p=0.049\),
on a statistic of \(0.004\).

Brier isolates external. Its mean absolute movement is \(0.023\), and the
sign-flip \(p\)-value against each of the other four blocks is \(0.008\).
The last-slice Brier interval is \([-0.044,-0.006]\). PO-risk does not
separate external from amount (mean absolute movement \(0.005\) and
\(0.004\), sign-flip \(p=0.253\)). On the last slice those two PO intervals
exclude zero, and the personal, region, and other intervals include zero.
HSIC deltas are all below \(0.001\).

Along application id the covariates do not separate. The peak-day default
model uses the external scores, so replacing that block moves the Brier
score transferred onto the later slices.
Tables~\ref{tab:home-credit-path}--\ref{tab:home-credit-sign} record the path.


\begin{table}[t]
\centering
\caption{Synthetic pathway. The anchor is day 0. Each step is one day closer to the drift peak. Shares are the block's fraction of the total absolute movement.}
\label{tab:pathway-synthetic}
\begin{tabular}{rcccccccc}
\toprule
step & PO-risk & Brier & MSE & driver PO & driver Brier & driver MSE & driver MMD & context MMD \\
\midrule
1 & 0.005 & 0.236 & 0.337 & 0.471 & 0.420 & 0.492 & 0.863 & 0.000 \\
2 & 0.014 & 0.231 & 0.403 & 0.791 & 0.672 & 0.548 & 0.886 & 0.076 \\
3 & 0.017 & 0.212 & 0.486 & 0.822 & 0.650 & 0.689 & 0.905 & 0.095 \\
4 & 0.026 & 0.195 & 0.901 & 0.821 & 0.868 & 0.810 & 0.920 & 0.071 \\
5 & 0.031 & 0.182 & 1.427 & 0.336 & 0.822 & 0.826 & 0.864 & 0.136 \\
6 & 0.032 & 0.158 & 2.589 & 0.923 & 0.911 & 0.921 & 0.852 & 0.144 \\
7 & 0.050 & 0.122 & 4.496 & 0.905 & 0.959 & 0.936 & 0.834 & 0.166 \\
\bottomrule
\end{tabular}
\end{table}

\begin{table}[t]
\centering
\caption{Synthetic day 7, the drift peak. Deltas are the change in the loss when that block is replaced by the TabDDPM draw from the anchor conditional. Intervals are bootstrap percentiles. MMD and HSIC $p$-values use 40 permutations.}
\label{tab:pathway-synthetic-day7}
\small
\resizebox{\textwidth}{!}{%
\begin{tabular}{lccccccccccc}
\toprule
block & PO $\Delta$ & PO interval & Brier $\Delta$ & Brier interval & MSE $\Delta$ & MMD & MMD $p$ & HSIC $\Delta$ & HSIC $p$ & PO share & Brier share \\
\midrule
stable & -0.002 & [-0.008, 0.004] & -0.004 & [-0.008, 0.001] & -0.292 & 0.000 & 1.000 & -0.0004 & 0.585 & 0.015 & 0.028 \\
context & -0.005 & [-0.016, 0.002] & -0.001 & [-0.002, -0.000] & -0.035 & 0.084 & 0.024 & -0.0002 & 0.341 & 0.041 & 0.010 \\
driver & -0.103 & [-0.122, -0.062] & -0.132 & [-0.139, -0.126] & -4.835 & 0.420 & 0.024 & 0.0003 & 0.049 & 0.905 & 0.959 \\
noise & 0.004 & [-0.002, 0.008] & -0.000 & [-0.002, 0.001] & 0.001 & 0.000 & 1.000 & -0.0002 & 0.024 & 0.039 & 0.003 \\
\bottomrule
\end{tabular}}
\end{table}

\begin{table}[t]
\centering
\caption{Sign-flip test over the seven synthetic steps. The contrast is the absolute movement of the driver minus the other block.}
\label{tab:pathway-synthetic-sign}
\begin{tabular}{llcc}
\toprule
loss & comparison & mean gap & $p$ \\
\midrule
PO-risk & driver $>$ stable & 0.031 & 0.008 \\
PO-risk & driver $>$ context & 0.029 & 0.008 \\
PO-risk & driver $>$ noise & 0.029 & 0.021 \\
Brier & driver $>$ stable & 0.050 & 0.014 \\
Brier & driver $>$ context & 0.060 & 0.008 \\
Brier & driver $>$ noise & 0.060 & 0.008 \\
MSE & driver $>$ stable & 1.511 & 0.014 \\
MSE & driver $>$ context & 1.821 & 0.008 \\
MSE & driver $>$ noise & 1.828 & 0.008 \\
\bottomrule
\end{tabular}
\end{table}

\begin{table}[t]
\centering
\caption{IEEE-CIS transaction day 180, seven days after the fraud-rate peak on day 173. Same columns as the synthetic peak day. No latent MSE on this table.}
\label{tab:pathway-ieee-day}
\small
\resizebox{\textwidth}{!}{%
\begin{tabular}{lcccccccccc}
\toprule
block & PO $\Delta$ & PO interval & Brier $\Delta$ & Brier interval & MMD & MMD $p$ & HSIC $\Delta$ & HSIC $p$ & PO share & Brier share \\
\midrule
amount & 0.001 & [-0.001, 0.006] & -0.000 & [-0.000, 0.000] & 0.002 & 0.098 & -0.0001 & 0.122 & 0.115 & 0.002 \\
card & -0.004 & [-0.014, 0.001] & -0.013 & [-0.023, -0.004] & 0.001 & 1.000 & -0.0003 & 0.098 & 0.324 & 0.420 \\
address & 0.006 & [-0.002, 0.013] & 0.005 & [0.002, 0.008] & 0.000 & 1.000 & 0.0000 & 0.073 & 0.484 & 0.159 \\
counts & 0.000 & [-0.006, 0.003] & -0.003 & [-0.006, 0.002] & 0.000 & 0.146 & -0.0000 & 0.122 & 0.007 & 0.102 \\
timedelta & 0.001 & [-0.007, 0.012] & 0.009 & [-0.003, 0.023] & 0.000 & 1.000 & -0.0001 & 0.049 & 0.061 & 0.302 \\
hour & 0.000 & [-0.003, 0.002] & -0.000 & [-0.005, 0.004] & 0.002 & 0.073 & 0.0002 & 0.024 & 0.010 & 0.015 \\
\bottomrule
\end{tabular}}
\end{table}

\begin{table}[t]
\centering
\caption{Sign-flip test over the seven IEEE-CIS steps after the peak. PO-risk does not name a single block. The transferred Brier movement is largest for the card block.}
\label{tab:pathway-ieee-sign}
\begin{tabular}{llcc}
\toprule
loss & comparison & mean gap & $p$ \\
\midrule
PO-risk & timedelta $>$ amount & 0.002 & 0.021 \\
PO-risk & timedelta $>$ card & 0.001 & 0.303 \\
PO-risk & timedelta $>$ address & 0.002 & 0.209 \\
PO-risk & timedelta $>$ counts & 0.002 & 0.178 \\
PO-risk & timedelta $>$ hour & 0.002 & 0.042 \\
Brier & card $>$ amount & 0.008 & 0.008 \\
Brier & card $>$ address & 0.006 & 0.008 \\
Brier & card $>$ counts & 0.005 & 0.028 \\
Brier & card $>$ timedelta & 0.002 & 0.292 \\
Brier & card $>$ hour & 0.004 & 0.054 \\
\bottomrule
\end{tabular}
\end{table}

\begin{table}[t]
\centering
\caption{MovieLens-1M, seven calendar days after a peak in the high-rating rate. Rating is the outcome. Shares use absolute movement.}
\label{tab:movielens-path}
\small
\resizebox{\textwidth}{!}{%
\begin{tabular}{lccccccccccc}
\toprule
step & rate & PO-risk & Brier & PO top & PO share & Brier top & Brier share & MSE top & MSE share & MMD top & MMD share \\
\midrule
2001-01-02 & 0.503 & 0.041 & 0.263 & user & 0.614 & context & 0.692 & user & 0.769 & user & 0.825 \\
2001-01-03 & 0.597 & 0.015 & 0.255 & item & 0.338 & item & 0.737 & user & 0.474 & user & 0.655 \\
2001-01-04 & 0.500 & 0.024 & 0.244 & item & 0.629 & context & 0.426 & item & 0.601 & context & 0.698 \\
2001-01-05 & 0.539 & 0.032 & 0.255 & item & 0.549 & item & 0.628 & item & 0.708 & user & 0.670 \\
2001-01-06 & 0.657 & 0.046 & 0.251 & user & 0.802 & item & 0.617 & item & 0.599 & context & 0.498 \\
2001-01-07 & 0.551 & 0.031 & 0.251 & item & 0.544 & item & 0.695 & item & 0.508 & context & 0.618 \\
2001-01-08 & 0.490 & 0.027 & 0.246 & user & 0.894 & item & 0.593 & item & 0.497 & context & 0.582 \\
\bottomrule
\end{tabular}}
\end{table}

\begin{table}[t]
\centering
\caption{MovieLens last day of the pathway. Rating MSE is the predictive companion of Brier.}
\label{tab:movielens-day}
\small
\resizebox{\textwidth}{!}{%
\begin{tabular}{lccccccccccc}
\toprule
block & PO $\Delta$ & PO interval & Brier $\Delta$ & Brier interval & MSE $\Delta$ & MMD & MMD $p$ & HSIC $\Delta$ & HSIC $p$ & PO share & Brier share \\
\midrule
user & -0.026 & [-0.048, -0.008] & -0.008 & [-0.014, 0.001] & -0.047 & 0.052 & 0.024 & 0.0021 & 0.024 & 0.894 & 0.284 \\
item & -0.000 & [-0.006, 0.004] & -0.018 & [-0.025, -0.010] & -0.048 & 0.006 & 0.024 & -0.0000 & 0.024 & 0.006 & 0.593 \\
context & 0.003 & [-0.001, 0.004] & 0.004 & [0.001, 0.006] & -0.002 & 0.081 & 0.024 & 0.0019 & 0.024 & 0.101 & 0.123 \\
\bottomrule
\end{tabular}}
\end{table}

\begin{table}[t]
\centering
\caption{Sign-flip test across the seven steps. The contrast is absolute movement.}
\label{tab:movielens-sign}
\begin{tabular}{llcc}
\toprule
loss & comparison & mean gap & $p$ \\
\midrule
PO-risk & user $>$ item & 0.004 & 0.198 \\
PO-risk & user $>$ context & 0.007 & 0.016 \\
Brier & item $>$ user & 0.005 & 0.029 \\
Brier & item $>$ context & 0.005 & 0.074 \\
MSE & item $>$ user & 0.017 & 0.125 \\
MSE & item $>$ context & 0.030 & 0.016 \\
\bottomrule
\end{tabular}
\end{table}

\begin{table}[t]
\centering
\caption{Home Credit along SK\_ID\_CURR. Seven slices after a peak in the default rate. There is no event time. Every block MMD stays below 0.005. An MMD entry of -- is a total below $10^{-4}$, and a large MMD share is still a split of a near-zero total.}
\label{tab:home-credit-path}
\small
\resizebox{\textwidth}{!}{%
\begin{tabular}{lccccccccc}
\toprule
step & rate & PO-risk & Brier & PO top & PO share & Brier top & Brier share & MMD top & MMD share \\
\midrule
slice 32 & 0.079 & 0.030 & 0.202 & other & 0.358 & external & 0.695 & external & 0.925 \\
slice 33 & 0.097 & 0.036 & 0.201 & external & 0.419 & external & 0.500 & other & 0.531 \\
slice 34 & 0.067 & 0.046 & 0.210 & other & 0.417 & external & 0.437 & amount & 0.423 \\
slice 35 & 0.084 & 0.034 & 0.193 & amount & 0.514 & external & 0.702 & external & 0.939 \\
slice 36 & 0.090 & 0.032 & 0.213 & external & 0.370 & external & 0.379 & -- & -- \\
slice 37 & 0.086 & 0.046 & 0.201 & external & 0.401 & external & 0.366 & -- & -- \\
slice 38 & 0.074 & 0.043 & 0.207 & amount & 0.386 & external & 0.508 & -- & -- \\
\bottomrule
\end{tabular}}
\end{table}

\begin{table}[t]
\centering
\caption{Home Credit last slice. PO-risk and transferred Brier on the default label.}
\label{tab:home-credit-day}
\small
\resizebox{\textwidth}{!}{%
\begin{tabular}{lcccccccccc}
\toprule
block & PO $\Delta$ & PO interval & Brier $\Delta$ & Brier interval & MMD & MMD $p$ & HSIC $\Delta$ & HSIC $p$ & PO share & Brier share \\
\midrule
external & -0.009 & [-0.022, -0.001] & -0.026 & [-0.044, -0.006] & 0.000 & 0.317 & -0.0005 & 0.024 & 0.362 & 0.508 \\
amount & -0.010 & [-0.022, -0.003] & -0.003 & [-0.006, 0.002] & 0.000 & 0.317 & -0.0001 & 0.512 & 0.386 & 0.058 \\
personal & -0.002 & [-0.009, 0.007] & 0.004 & [-0.002, 0.010] & 0.000 & 1.000 & -0.0000 & 0.195 & 0.071 & 0.076 \\
region & -0.003 & [-0.005, 0.005] & 0.015 & [0.010, 0.024] & 0.000 & 0.195 & -0.0002 & 0.390 & 0.118 & 0.284 \\
other & 0.002 & [-0.009, 0.008] & 0.004 & [-0.006, 0.010] & 0.000 & 0.049 & 0.0001 & 0.073 & 0.064 & 0.074 \\
\bottomrule
\end{tabular}}
\end{table}

\begin{table}[t]
\centering
\caption{Sign-flip test across the seven steps. The contrast is absolute movement.}
\label{tab:home-credit-sign}
\begin{tabular}{llcc}
\toprule
loss & comparison & mean gap & $p$ \\
\midrule
PO-risk & external $>$ amount & 0.001 & 0.253 \\
PO-risk & external $>$ personal & 0.004 & 0.008 \\
PO-risk & external $>$ region & 0.003 & 0.023 \\
PO-risk & external $>$ other & 0.002 & 0.078 \\
Brier & external $>$ amount & 0.020 & 0.008 \\
Brier & external $>$ personal & 0.018 & 0.008 \\
Brier & external $>$ region & 0.015 & 0.008 \\
Brier & external $>$ other & 0.018 & 0.008 \\
\bottomrule
\end{tabular}
\end{table}


\end{document}
